%--Packages
\documentclass{article}
\usepackage[utf8]{inputenc}
\usepackage{titling}
\usepackage{fancyhdr} %--for fancy header/ footers
\usepackage{tikz} %--for border in title page
\usetikzlibrary{calc} %-- access to tikslibrary commands
\usepackage{hyperref}

%--Parameters
\title{Elaboration Exercise I}
\author{Aaron Hammond\\43691455}
\date{August 2019}
\fancyhfoffset[]{}
\fancy

%--Page Setup
\pagestyle{fancy}
\fancyhf{}

\rhead{Scoping Exercise II}
\lhead{Digital Humanities}

\fancyfoot[RE,RO]{Aaron Hammond | 43691455}
\fancyfoot[LE,LO]{Page: \thepage}

\renewcommand{\headrulewidth}{2pt}
\renewcommand{\footrulewidth}{1pt}
%--Commands

\newcommand\HRule{\rule{\textwidth}{1pt}} %-- horizontal lines for formatting

%--
\begin{document}

%--Title page

\begin{titlepage}

\begin{tikzpicture}[remember picture, overlay]
  \draw[line width = 4pt] ($(current page.north west) + (1in,-1in)$) rectangle ($(current page.south east) + (-1in,1in)$);
\end{tikzpicture}

\begin{center}
\HRule\\[0.4cm]
\huge{Elaboration Document I}\\[0.5cm]
\huge{Aaron Hammond}\\
\large{\today}\\[0.4cm]

\HRule \\[1cm]


\noindent This document is to elaborate on my proof of concept's scoping document and begin to think about the practical ways of achieving each  of the following steps:
\end{center}

\def\contentsname{\empty} %--Remove contents header from toc
\tableofcontents

\end{titlepage}

%--

\section{Transcribe audio into plain text}
\section{Convert articles into rich text}
\section{Create plain text files from field notes}
\section{Add headers to all library files}
\section{Analyze files and produce report}
\

\end{document}
